%% Title
\titlepage[of Emmanuel College]{%
\vspace{-4cm}
\centering
\includegraphics[width=0.2\textwidth]{Miscellaneous/Emmanuel_College_Crest.png} \\
A dissertation submitted to the University of Cambridge\\ for the degree of Doctor of Philosophy \\ August 2017}

%% Abstract
\begin{abstract}[\smaller \thetitle\\ \vspace*{1cm} \smaller {\theauthor}]
  %\thispagestyle{empty}
This thesis describes the optimisation of the calorimeter design for collider experiments at the future Compact LInear Collider (CLIC) and the International Linear Collider (ILC).  The detector design of these experiments is built around high-granularity Particle Flow Calorimetry that, in contrast to traditional calorimetry, uses the energy measurements for charged particles from the tracking detectors. This can only be realised if calorimetric energy deposits from charged particles can be separated from those of neutral particles. This is made possible with fine granularity calorimeters and sophisticated pattern recognition software, which is provided by the PandoraPFA algorithm. This thesis presents results on Particle Flow calorimetry performance for a number of detector configurations. To obtain these results a new calibration procedure was developed and applied to the detector simulation and reconstruction to ensure optimal performance was achieved for each detector configuration considered.

This thesis also describes the development of a software compensation technique that vastly improves the intrinsic energy resolution of a Particle Flow Calorimetry detector. This technique is implemented within the PandoraPFA framework and demonstrates the gains that can be made by fully exploiting the information provided by the fine granularity calorimeters envisaged at a future linear collider.

A study of the sensitivity of the CLIC experiment to anomalous gauge couplings that {affect} vector boson scattering processes is presented. These anomalous couplings provide insight into possible beyond standard model physics. This study, which utilises the excellent jet energy resolution from Particle Flow Calorimetry, was performed at centre-of-mass energies of 1.4 TeV and 3 TeV with integrated luminosities of 1.5$\text{ab}^{-1}$ and 2$\text{ab}^{-1}$ respectively. The precision achievable at CLIC is shown to be approximately one to two orders of magnitude better than that currently offered by the LHC.

In addition, a study into various technology options for the CLIC vertex detector is described.
\end{abstract}

%% Declaration
\begin{declaration}
This dissertation is the result of my own work, except where explicit reference is made to the work of others, and has not been submitted for another qualification to this or any other university. This dissertation does not exceed the word limit for the respective Degree Committee.
  \vspace*{1cm}
  \begin{flushright}
    Steven Green
  \end{flushright}
\end{declaration}

%% Acknowledgements
\begin{acknowledgements}
Firstly, I would like to thank Mark Thomson and John Marshall for the excellent supervision that they have both given me throughout the course of this work.  They have both been sources of insight, wisdom and encouragement over the past years and I am truly grateful to have had them as my supervisors.  I would like to thank Boruo Xu for his friendship and sense of humour that has made sharing an office with him for the past four years a joy.  I am extremely grateful to have worked in the Cambridge HEP group throughout my PhD and would like to thank everyone there for creating a welcoming and friendly working environment.  

The work presented in this thesis would not have been possible without the support of the CLICdp and ILD collaborations.  I am very thankful to the CLICdp group for supporting me during my time at CERN.  In particular, I would like to thank Lucie Linssen, Dominik Dannheim and Daniel Hynds for their support and guidance on the CLICdp vertex research in this thesis.  Funding for this research was provided by the UK STFC and I would like to thank them for their generous support. 

I would like to thank my parents, family and friends for all of their help getting me to this point.  The foundation that my parents provided has given me the platform to get to where I am now and I am eternally grateful for all that they have done for me.  

I would like to give special thanks to my wife Emily.  She has been an inspirational support to me even since we first met and I lack the words to truly describe how thankful I am for her.  Her positivity and cheerful nature have been invaluable to me personally, while her checking of gammer has been equally invaluable to the reader.  

Finally, I would like to thank God for graciously giving me the opportunity to try and figure out how He put the universe together.  

\end{acknowledgements}

%% Preface
%\begin{preface}
%  This thesis describes my research on various aspects of the \LHCb
%  particle physics program, centred around the \LHCb detector and \LHC
%  accelerator at \CERN in Geneva.

%  \noindent
%  For this example, I'll just mention \ChapterRef{chap:SomeStuff}
%  and \ChapterRef{chap:MoreStuff}.
%\end{preface}

%% ToC
\tableofcontents

%% Strictly optional!
\frontquote{%
I thank and praise you, God of my ancestors: You have given me wisdom and power}%
{Daniel 2:23a}
%% I don't want a page number on the following blank page either.
\thispagestyle{empty}
